%% kapa: exam=1 round=31 year=2020 period=1 form=A subject=real_estate subject_ko=부동산학원론 seq=30 number=110
%% source: Q-Net 감정평가사 기출문제 (cst00309) artlSeq=5207219 file=extracted/1008/5207219/A형/A형/2020년도 제31회 감정평가사 1차 시험_1교시_A형.pdf
%% conversion: kordoc 4.24.51 -> tools/gen1.py (auto); review=proofread-sonnet; flags=none
\begin{question}{110}
화폐의 시간가치에 관한 설명으로 옳지 않은 것은? (단, 다른 조건은 동일함)
\begin{choices}
  \item 은행으로부터 주택구입자금을 원리금균등분할상환 방식으로 대출한 가구가 매월 상환할 원리금을 계산하는 경우, 저당상수를 사용한다.
  \item 일시불의 미래가치계수는 이자율이 상승할수록 커진다.
  \item 연금의 현재가치계수와 저당상수는 역수관계이다.
  \item 연금의 미래가치계수와 감채기금계수는 역수관계이다.
  \item 3년 후에 주택자금 5억원을 만들기 위해 매 기간 납입해야 할 금액을 계산하는 경우, 연금의 미래가치계수를 사용한다.
\end{choices}
\end{question}
\answer{5}
